\section{Conclusion and Future Work}
\label{sec:conclusion}

This paper set out to test whether standard imbalance corrections
remain worth applying once the labels they depend on cannot be
trusted. Across 26 datasets, seven correction methods spanning two
model families, and two directed label-error mechanisms, every
correction audited underperformed its matched uncorrected baseline in
at least some tested condition. How often and how badly it did so
depended on the method, on the direction of the label error, and on
whether balanced accuracy or calibration was the metric being judged,
never on a single number alone.
The simplest correction tested, a threshold shift keyed to the class
prior, was consistently the most robust on discrimination, leading
every other method on Correction Harm Rate and mean Correction Gain in
both directions. The two most mathematically involved corrections,
DistPFN and DistPFN-T, held the highest Correction Harm Rate of any
method under both directions, more than five times the next-highest
method's rate under Direction B, the opposite of what their added
sophistication might suggest, yet those same two corrections were also
the best-calibrated methods in the entire study. A robustness check
restricted to genuinely imbalanced training sets, excluding the
balanced $r{=}0.50$ control, shows a real share, though not the
majority, of that harm rate concentrated at the balanced condition
specifically: the DistPFN pair's Direction B Correction Harm Rate
falls from 43.8\%/47.4\% to 29.8\%/37.8\% once $r{=}0.50$ is excluded,
while \texttt{tabpfn\_threshold}'s lead over the rest of the field
holds or strengthens throughout. A correction that more frequently
degrades discrimination relative to its uncorrected baseline can
nevertheless produce better-calibrated probabilities, a split this
paper's results surface rather than resolve in either method's favor.

The five contribution metrics introduced here, Correction Gain,
Correction Harm Rate, Regret, the break-even contamination rate, and
the Prior Sensitivity Gap, made this possible, separating whether a
correction helps from how often it fails, how badly it fails when it
does, and how much of that failure traces to the corrupted prior
specifically. The same framework also
surfaced a result this paper did not anticipate, that the corrupted
prior is not the primary reason TabPFN's posterior-rescaling
corrections decline under contamination, since handing them the true
prior did not reliably help, and on average even hurt slightly.
Section~\ref{sec:discussion-prior} offers one explanation, that these
corrections never touch the mislabeled rows still sitting in TabPFN's
in-context prompt.

Several directions follow from what this study leaves open. Confirming
the explanation above would need an experiment that cleans the
mislabeled rows from TabPFN's in-context prompt directly, rather than
only adjusting the prior applied to it; separately, because training
here is deliberately pushed to each tested ratio while the test
partition keeps each dataset's natural class distribution, a complementary
matched-prior evaluation regime, resampling the test partition to match
each tested $r$ rather than keeping it at the dataset's natural
prevalence, would isolate how much of a prior-adjusting method's
robustness, particularly \texttt{tabpfn\_threshold}'s, reflects imbalance
tolerance specifically rather than the train/test prior mismatch the
present design creates; it targets a different, less deployment-oriented
question than the one this study poses, so it belongs alongside these
results as a complementary direction rather than as an outstanding gap in
them. Beyond those two, testing whether
the same robustness ranking holds on a dataset whose minority class
alone permits a larger controlled training set, on tabular foundation
models other than TabPFN, under multi-class imbalance, and under
instance-dependent or simultaneous-direction label noise rather than
the class-conditional, single-direction mechanisms audited here, are
the most direct ways to extend this audit.